\documentclass[12pt]{article}
\usepackage[letterpaper,top=0.75in,bottom=0.75in,left=0.78in,right=1.55in,marginparwidth=1.15in,marginparsep=0.18in]{geometry}
\usepackage{fontspec}
\setmainfont{Latin Modern Roman}
\usepackage{microtype}
\usepackage{fancyhdr}
\usepackage{titlesec}
\usepackage{enumitem}
\usepackage{xcolor}
\usepackage[hidelinks]{hyperref}
\setlength{\parindent}{1.05em}
\setlength{\parskip}{0.32em}
\setlength{\headheight}{15pt}
\titleformat{\section}{\large\bfseries\scshape}{}{0pt}{}
\titleformat{\subsection}{\normalsize\bfseries}{}{0pt}{}
\pagestyle{fancy}
\fancyhf{}
\fancyhead[L]{Whately Logic: Macbeth Lit-Anchor}
\fancyhead[R]{Lesson 3}
\fancyfoot[C]{\thepage}
\renewcommand{\headrulewidth}{0.4pt}

\newcommand{\focusbox}[1]{\par\vspace{0.35em}\noindent\begingroup\setlength{\fboxsep}{6pt}\colorbox{gray!12}{\begin{minipage}{\dimexpr\linewidth-2\fboxsep\relax}#1\end{minipage}}\endgroup\par\vspace{0.45em}}
\newcommand{\studentline}{\par\vspace{0.68em}\hrule\vspace{0.72em}}
\newcommand{\gloss}[1]{\marginpar{\scriptsize\raggedright\setlength{\baselineskip}{7.8pt}\color{gray!70!black}#1}}
\newcommand{\speaker}[1]{\par\vspace{0.35em}\noindent\textsc{#1}\par\vspace{-0.15em}}

\begin{document}
\begin{center}
{\LARGE\bfseries Macbeth Lit-Anchor}\par\vspace{0.18em}
{\large\scshape What Does ``A Man'' Mean?}\par\vspace{0.35em}
{Lesson 3: Propositions, Quality, Quantity, and Clear Claims}\par\vspace{0.55em}
\rule{0.82\textwidth}{0.4pt}
\end{center}

\section*{Purpose}
A lit-anchor connects the week's logic lesson to a recurring literary text. In Lesson 3, Whately asks us to look closely at propositions: statements that can be true or false. Before we can test an argument, we must know what is being asserted, what the subject is, what the predicate is, whether the claim affirms or denies something, and whether it speaks about all, none, or some members of a class.

\textit{Macbeth} gives us a sharp example. Lady Macbeth does not merely insult Macbeth. She pressures him with claims about what a man is. Macbeth answers with a rival claim: daring can go beyond what truly belongs to manhood. Their quarrel is dramatic, but it is also logical. The first task is not to decide who sounds stronger. The first task is to state the hidden propositions clearly.

\focusbox{\textbf{Whately's Lesson 3 habit:} Turn ordinary language into clear proposition form. Identify the subject, predicate, quality, and quantity before judging whether the claim is true.}

\section*{Central Question}
When Macbeth and Lady Macbeth argue about being ``a man,'' what exact propositions are they making about courage, murder, and moral limits?

\section*{Guided Text: Act 1, Scene 7}
Read the original text below. The side notes are light glosses only; they do not replace the text. Watch for the way Lady Macbeth turns courage and manhood into a test, and watch how Macbeth tries to define manhood differently.

\speaker{Lady Macbeth}
\begin{verse}
Was the hope drunk\\
Wherein you dress'd yourself? hath it slept since?\\
And wakes it now, to look so green and pale\\
At what it did so freely? From this time\\
Such I account thy love. Art thou afeard\\
To be the same in thine own act and valour\\
As thou art in desire? Wouldst thou have that\\
Which thou esteem'st the ornament of life,\\
And live a coward in thine own esteem,\\
Letting ``I dare not'' wait upon ``I would,''\\
Like the poor cat i' the adage?\gloss{Lady Macbeth frames hesitation as cowardice: Macbeth wants the crown but says, ``I dare not.''}
\end{verse}

\speaker{Macbeth}
\begin{verse}
Prithee, peace:\\
I dare do all that may become a man;\\
Who dares do more is none.\gloss{Macbeth gives a definition: true manhood has moral limits. More daring is not always more manly.}
\end{verse}

\speaker{Lady Macbeth}
\begin{verse}
What beast was't, then,\\
That made you break this enterprise to me?\\
When you durst do it, then you were a man;\\
And, to be more than what you were, you would\\
Be so much more the man. Nor time nor place\\
Did then adhere, and yet you would make both:\\
They have made themselves, and that their fitness now\\
Does unmake you.\gloss{She answers by making daring the measure of manhood: when Macbeth dared the murder, he was a man. More daring would make him more manly.}
\end{verse}

\speaker{Macbeth}
\begin{verse}
If we should fail?\gloss{Macbeth's question turns from moral definition to practical risk.}
\end{verse}

\speaker{Lady Macbeth}
\begin{verse}
We fail!\\
But screw your courage to the sticking-place,\\
And we'll not fail.\gloss{Lady Macbeth shifts the focus from whether the act is right to whether they can succeed.}
\end{verse}

\section*{Guided Analysis}
Lady Macbeth's language is powerful because it sounds personal and emotional. Whately helps us slow it down. A taunt can hide a proposition. Once the proposition is clear, it can be tested.

\subsection*{Claim 1: Lady Macbeth's implied test of manhood}
\begin{itemize}[leftmargin=*]
\item \textbf{Ordinary wording:} ``When you durst do it, then you were a man.''
\item \textbf{Possible standard form:} All truly manly people are people who dare this deed.
\item \textbf{Subject:} truly manly people.
\item \textbf{Predicate:} people who dare this deed.
\item \textbf{Quality:} affirmative.
\item \textbf{Quantity:} universal.
\item \textbf{Problem exposed by clear form:} Lady Macbeth makes murder sound like the test of manhood. Once the claim is written clearly, we can ask whether that predicate belongs in the definition of a man at all.
\end{itemize}

\subsection*{Claim 2: Macbeth's moral boundary}
\begin{itemize}[leftmargin=*]
\item \textbf{Ordinary wording:} ``I dare do all that may become a man; / Who dares do more is none.''
\item \textbf{Possible standard form:} No people who dare beyond what befits a man are true men.
\item \textbf{Subject:} people who dare beyond what befits a man.
\item \textbf{Predicate:} true men.
\item \textbf{Quality:} negative.
\item \textbf{Quantity:} universal.
\item \textbf{Problem exposed by clear form:} Macbeth does not deny courage. He denies that courage without moral limit is true manhood.
\end{itemize}

\subsection*{Claim 3: A clearer counter-proposition}
A careful reader might answer Lady Macbeth with an even clearer proposition:

\begin{itemize}[leftmargin=*]
\item \textbf{Standard form:} Some daring acts are not honorable acts.
\item \textbf{Form:} O proposition: Some S are not P.
\item \textbf{Subject:} daring acts.
\item \textbf{Predicate:} honorable acts.
\item \textbf{Why it matters:} This claim separates courage from goodness. It allows a person to be bold and still wrong.
\end{itemize}

\focusbox{\textbf{Lesson 3 takeaway:} Macbeth and Lady Macbeth disagree partly because they are using dramatic language instead of clear propositions. Whately's method exposes the hidden claim: Lady Macbeth treats murder as proof of manhood, while Macbeth briefly treats moral limit as part of manhood.}

\section*{Discussion}
Use these questions to prepare for the independent passage.

\begin{enumerate}[leftmargin=*]
\item What subject is Lady Macbeth really talking about when she says, ``then you were a man''?
\item What predicate does she attach to that subject?
\item Is her implied claim universal or particular? How can you tell?
\item Macbeth says, ``Who dares do more is none.'' What kind of proposition does that become when stated clearly?
\item Why is ``Some daring acts are not honorable acts'' a useful proposition for analyzing this scene?
\item How does clear proposition form make Lady Macbeth's pressure easier to challenge?
\end{enumerate}

\section*{Independent Text: Act 3, Scene 4}
Now read a later scene. Macbeth has murdered Duncan and arranged Banquo's murder. At the banquet, he sees Banquo's ghost. Lady Macbeth again questions his manhood. This time, notice whether Macbeth's idea of manhood has changed.

\speaker{Lady Macbeth}
\begin{verse}
Are you a man?\gloss{The question is short, but it still pressures Macbeth with a standard of manhood.}
\end{verse}

\speaker{Macbeth}
\begin{verse}
Ay, and a bold one, that dare look on that\\
Which might appal the devil.\gloss{Macbeth defines himself by boldness: he dares to look at what would terrify even the devil.}
\end{verse}

\speaker{Lady Macbeth}
\begin{verse}
O proper stuff!\\
This is the very painting of your fear:\\
This is the air-drawn dagger which, you said,\\
Led you to Duncan. O, these flaws and starts,\\
Impostors to true fear, would well become\\
A woman's story at a winter's fire,\\
Authorized by her grandam. Shame itself!\\
Why do you make such faces? When all's done,\\
You look but on a stool.\gloss{She treats his fear as shameful and unreal, as if no real threat is present.}
\end{verse}

\speaker{Macbeth}
\begin{verse}
What man dare, I dare:\\
Approach thou like the rugged Russian bear,\\
The arm'd rhinoceros, or the Hyrcan tiger;\\
Take any shape but that, and my firm nerves\\
Shall never tremble: or be alive again,\\
And dare me to the desert with thy sword;\\
If trembling I inhabit then, protest me\\
The baby of a girl. Hence, horrible shadow!\\
Unreal mockery, hence!\gloss{Macbeth claims universal courage: whatever any man dares, he dares. But the ghost still breaks him.}
\end{verse}

\speaker{Macbeth}
\begin{verse}
Why, so: being gone,\\
I am a man again.\gloss{Macbeth links recovered composure with being a man again.}
\end{verse}

\section*{Independent Analysis}
Answer in complete thoughts. Leave yourself enough evidence that this work can become a portfolio artifact.

\begin{enumerate}[leftmargin=*]
\item Translate Lady Macbeth's question, ``Are you a man?'' into one possible proposition. What is the subject? What is the predicate?\studentline\studentline
\item Translate Macbeth's claim, ``What man dare, I dare,'' into standard form. Is it universal or particular? Affirmative or negative?\studentline\studentline\studentline
\item Does Macbeth now define manhood by moral limits, by boldness, by lack of fear, or by something else? Use one line as evidence.\studentline\studentline\studentline
\item Write one A-form proposition about Macbeth in this scene: ``All S are P.''\studentline\studentline
\item Write one O-form proposition about Macbeth in this scene: ``Some S are not P.''\studentline\studentline
\item Which proposition is easier to defend from the text: ``All brave men are morally steady'' or ``Some bold men are morally disordered''? Explain.\studentline\studentline\studentline
\end{enumerate}

\section*{Portfolio Artifact: Macbeth Proposition Translation Chart}
Choose one line from this lit-anchor, or another claim about Macbeth's manhood, courage, fear, or guilt. Then complete the chart below.

\textbf{Original wording:}\studentline
\textbf{Ordinary-language claim:}\studentline\studentline
\textbf{Standard logical form:}\studentline
\textbf{Subject:}\studentline
\textbf{Predicate:}\studentline
\textbf{Quality:}\studentline
\textbf{Quantity:}\studentline
\textbf{What becomes clearer once the proposition is stated precisely?}\studentline\studentline\studentline

\section*{Closing Synthesis}
Whately teaches that clear reasoning begins with clear propositions. Macbeth shows why this matters in literature and in life. A vague challenge like ``be a man'' can feel powerful because its exact claim is hidden. Once the claim is translated, the pressure becomes testable. If Lady Macbeth's proposition is ``All true men dare this murder,'' Macbeth can answer: ``Some daring acts are not honorable acts.'' The tragedy is that Macbeth briefly sees the clearer proposition, then lets ambition, fear, and shame pull him away from it.

\end{document}
